\chapter{BusyBox}
\label{ch:busybox}

BusyBox is the Swiss Army knife of embedded Linux — a single binary that
provides over 300 Unix utilities. It is the foundation of Veilbox's
userspace, providing the shell, file utilities, init system, and network
tools.

\section{What is BusyBox?}

BusyBox was created in 1999 by Bruce Perens to build a complete Linux
system on a single floppy disk. The concept is simple: a single
statically-linked binary contains multiple ``applets,'' each implementing
a standard Unix utility.

When invoked through a symlink (e.g., \texttt{/bin/ls} $\rightarrow$
\texttt{/bin/busybox}), BusyBox checks \texttt{argv[0]} to determine
which applet to run:

\begin{lstlisting}
int main(int argc, char *argv[]) {
    const char *applet = bb_basename(argv[0]);
    // Strip leading "busybox" if it's a direct call
    // Lookup applet in applet table
    // Jump to applet's main function
}
\end{lstlisting}

\section{Applet Dispatch}

BusyBox's \texttt{main()} function uses a binary search through a sorted
table of all applet names. Each entry in the table is a triple:

\begin{lstlisting}
struct busybox_applet {
    const char *name;       // Applet name (e.g., "ls")
    int (*main)(int, char**);  // Function pointer
    unsigned flags;          // Feature flags
};
\end{lstlisting}

The table is sorted alphabetically at build time, enabling O(log n) lookup.

\section{BusyBox Configuration}

Like the Linux kernel, BusyBox uses Kconfig for configuration:

\begin{lstlisting}[language=bash]
# Configure BusyBox
make -C busybox-1.35.0 defconfig

# Enable specific applets
make -C busybox-1.35.0 menuconfig

# Build
make -C busybox-1.35.0 -j$(nproc)

# Install to directory
make -C busybox-1.35.0 CONFIG_PREFIX=/path/to/rootfs install
\end{lstlisting}

\section{Key Applets in Veilbox}

The following BusyBox applets are essential to Veilbox's operation:

\subsection{Shell: ash}

The Almquist Shell (ash) is BusyBox's default shell. It is a lightweight
POSIX-compatible shell derived from the NetBSD sh:

\begin{itemize}[nosep]
  \item \textbf{Size}: Approximately 80 KB (vs. bash at 900 KB).
  \item \textbf{Features}: Job control, functions, aliases, environment
        variables, I/O redirection, pipes, here-docs.
  \item \textbf{Limitations}: No arrays, no associative arrays, no
        process substitution (but these are bashisms).
\end{itemize}

\subsection{Init: /sbin/init}

BusyBox implements a minimal init system that reads /etc/inittab:

\begin{lstlisting}
# /etc/inittab
::sysinit:/etc/init.d/rcS
tty1::respawn:/sbin/getty 38400 tty1
ttyS0::respawn:/sbin/getty 38400 ttyS0
::shutdown:/bin/umount -a -r
\end{lstlisting}

Actions:
\begin{itemize}[nosep]
  \item \texttt{sysinit}: Run once at boot, before any respawn.
  \item \texttt{respawn}: Restart the process whenever it exits.
  \item \texttt{askfirst}: Like respawn but prompts before starting.
  \item \texttt{shutdown}: Run on SIGTERM/SIGUSR1.
  \item \texttt{restart}: Run when init receives SIGHUP.
\end{itemize}

\subsection{Core Utilities}

\begin{longtable}{|p{3cm}|p{4cm}|p{5cm}|}
\hline
\textbf{Applet} & \textbf{Purpose} & \textbf{Common Options} \\
\hline
ls & List directory contents & -l, -a, -h, -R \\
cp & Copy files & -r, -p, -f, -a \\
mv & Move/rename files & -f, -i, -n \\
rm & Remove files & -r, -f, -v \\
mkdir & Create directory & -p (parents) \\
cat & Concatenate files & -n, -b \\
echo & Print text & -e (escape), -n \\
grep & Search text & -r, -i, -v, -l \\
sed & Stream editor & -i, -n, -e \\
awk & Pattern scanning & (BusyBox awk, limited) \\
\hline
\end{longtable}

\subsection{Process Utilities}

\begin{lstlisting}
ps          # List processes
top         # Process monitor (simplified)
kill        # Send signal
pgrep       # Find process by name
pkill       # Kill process by name
\end{lstlisting}

\subsection{Network Utilities}

\begin{lstlisting}
ifconfig    # Configure network interface
route       # Show/manipulate routing table
ping        # ICMP echo
udhcpc      # DHCP client
netstat     # Network statistics
wget        # HTTP/FTP client
\end{lstlisting}

\subsection{Editors and Utilities}

\begin{lstlisting}
vi          # Minimal vi editor
less        # Pager (simpler than full less)
tail        # Show last N lines
head        # Show first N lines
sort        # Sort lines
uniq        # Remove duplicates
tee         # Split output to file and stdout
\end{lstlisting}

\section{Symbolic Link Installation}

The \texttt{make install} step creates symbolic links for each enabled
applet. For example:

\begin{lstlisting}
/bin/busybox -> /bin/busybox  (the binary itself)
/bin/sh      -> /bin/busybox
/bin/ls      -> /bin/busybox
/bin/cp      -> /bin/busybox
/usr/bin/vi  -> /bin/busybox
/sbin/init   -> /bin/busybox
\end{lstlisting}

If the \texttt{CONFIG\_PREFIX} is set to \texttt{rootfs/}, the symlinks
are created under \texttt{rootfs/bin/}, \texttt{rootfs/sbin/}, etc.

\section{Trade-offs}

BusyBox applets are not full replacements for their GNU counterparts:

\begin{itemize}[nosep]
  \item \textbf{Performance}: BusyBox grep is slower on large files
        (no mmap optimization). BusyBox sort is slower on large datasets.
  \item \textbf{Features}: BusyBox vi lacks syntax highlighting and
        split windows. BusyBox top lacks interactive commands.
  \item \textbf{Compatibility}: Some GNU-specific options are missing.
        Scripts that use \texttt{sed -r} need to use \texttt{-E} instead.
\end{itemize}

However, for a minimal container host these limitations rarely matter
--- containers typically provide their own userspace via images.

\section{Building BusyBox from Source}

\subsection{Download and Extract}

\begin{lstlisting}[language=bash]
curl -LO https://busybox.net/downloads/busybox-1.35.0.tar.bz2
tar xf busybox-1.35.0.tar.bz2
cd busybox-1.35.0
\end{lstlisting}

\subsection{Configuration}

BusyBox uses Kconfig like the Linux kernel:

\begin{lstlisting}[language=bash]
# Default configuration
make defconfig

# Customize (optional)
make menuconfig

# Key settings to verify:
# - Busybox Settings -> Build static binary (no shared libs)
# - Init Utilities -> init (enable)
# - Shells -> ash (enable)
# - Coreutils -> all standard tools
# - Network Utilities -> udhcpc, ifconfig, route, ping, wget
\end{lstlisting}

\subsection{Compilation}

\begin{lstlisting}[language=bash]
# Build (parallel)
make -j$(nproc)

# Install to rootfs directory
make CONFIG_PREFIX=/path/to/rootfs install

# Verify
ls -la /path/to/rootfs/bin/sh
/path/to/rootfs/bin/busybox --help
\end{lstlisting}

\subsection{Static vs Dynamic Linking}

\begin{itemize}[nosep]
  \item \textbf{Static}: \texttt{CONFIG\_STATIC=y}. Produces a ~900 KB
        binary with no external dependencies. Recommended for initramfs
        use. The binary is slightly larger because libc is included.
  \item \textbf{Dynamic}: \texttt{CONFIG\_STATIC=n}. Produces a ~300 KB
        binary but requires libc.so and ld-linux.so on the target.
        Slightly smaller overall but more complex to deploy.
\end{itemize}

Veilbox uses static linking for maximum portability and simplicity.

\section{The Ash Shell}

Ash is the Almquist Shell, a lightweight POSIX-compatible shell:

\begin{longtable}{|p{3cm}|p{4cm}|p{5cm}|}
\hline
\textbf{Feature} & \textbf{Ash} & \textbf{Bash} \\
\hline
Size & ~80 KB & ~900 KB \\
Arrays & No & Yes \\
Associative arrays & No & Yes \\
Process substitution & No & Yes \\
Here-strings & No & Yes \\
regex =~ operator & No & Yes \\
printf builtin & Yes & Yes \\
Job control & Yes & Yes \\
Functions & Yes & Yes \\
Local variables & Yes & Yes \\
\hline
\end{longtable}

\subsection{Ash Scripting Tips}

\begin{itemize}[nosep]
  \item Use \texttt{\$\{var:-default\}} for default values.
  \item Use \texttt{test} or \texttt{[ ]} for conditionals (not [[ ]]).
  \item Use \texttt{\$(cmd)} for command substitution (not backticks).
  \item Use \texttt{local var=value} for function-local variables.
  \item Use \texttt{read -r} to read lines without backslash escaping.
  \item Redirect stderr with \texttt{2>\&1}.
  \item Use \texttt{trap} for signal handling.
  \item Use \texttt{type command} to check if a command exists.
\end{itemize}
